% mechanism_test_methodology.tex
% Stand-alone methodology snippet for the mechanism (gap/surplus) test.
% Designed to be \input{} into the paper. No \documentclass.
% Required packages in main document: amsmath, amssymb, bm.

\subsection{Mechanism test: do treated firms learn from connected partners?}
\label{sec:mechanism_test}

If the spillover effect runs through CBA content diffusing along worker
flows, the same channel should also operate \emph{from} treated firms'
connected partners \emph{back} to the treated firm itself. Specifically,
when treated firms renegotiate post-Súmula 277, they should add the most
new clauses in the dimensions where their connected, untreated partners
already had a head start pre-reform.

We test this prediction directly on the treated balanced-panel sample,
reshaped long over the 139 clause types $c$.

\paragraph{Pre-treatment mismatch.} For each treated firm $i$ and each
clause type $c$, we summarize what its connected, untreated worker-flow
partners $j$ already carried in their CBAs pre-reform. Let $w_{ij}$ be the
average of the four pre-treatment year-pair flow ratios from $i$ to $j$
($\overline{w_{ij}} = \tfrac14(r_{ij}^{07{-}08} + r_{ij}^{08{-}09} + r_{ij}^{09{-}10} + r_{ij}^{10{-}11})$),
and let $c^{\text{pre}}_{jc}$ be partner $j$'s pre-treatment count of
clauses of type $c$. The connectivity-weighted partner average is
%
\begin{equation}
  \bar{c}_{-i,c}
  \;=\; \frac{\sum_{j \in \mathcal{P}_i} w_{ij}\, c^{\text{pre}}_{jc}}
             {\sum_{j \in \mathcal{P}_i} w_{ij}},
\end{equation}
%
where $\mathcal{P}_i$ is the set of untreated firms with positive worker
flow to $i$ during 2007--2011. Note that the weights normalize out: the
quantity is a content-distance summary, not an exposure dose. From this we
build three pre-treatment, time-invariant mismatch measures:
%
\begin{align}
  \text{gap}_{ic}     &\;=\; \max\!\bigl(0,\; \bar{c}_{-i,c} - c^{\text{pre}}_{ic}\bigr)
  & \text{(partners had more)}\\
  \text{surplus}_{ic} &\;=\; \max\!\bigl(0,\; c^{\text{pre}}_{ic} - \bar{c}_{-i,c}\bigr)
  & \text{(firm $i$ had more)}\\
  \text{raw\_gap}_{ic} &\;=\; \bar{c}_{-i,c} - c^{\text{pre}}_{ic}
  & \text{(signed)}.
\end{align}
%
By construction $\text{gap}_{ic}$ and $\text{surplus}_{ic}$ have disjoint
support at the firm--clause cell level (only one is positive in any cell,
the other is zero), and $\text{raw\_gap}_{ic} = \text{gap}_{ic} - \text{surplus}_{ic}$.

\paragraph{Specification.} On the long panel of treated balanced-panel
firms $\times$ years $\times$ clause types, we estimate
%
\begin{equation}
  cl\_count_{ict}
  \;=\; \beta \,\bigl(X_{ic} \times T_t\bigr)
  \;+\; \alpha_{ic}
  \;+\; \alpha_{ft}
  \;+\; \varepsilon_{ict},
  \label{eq:mechanism_spec}
\end{equation}
%
where $X_{ic} \in \{\text{gap}_{ic},\, \text{surplus}_{ic},\, \text{raw\_gap}_{ic}\}$,
$\alpha_{ic}$ is a firm $\times$ clause-type fixed effect (which absorbs
the time-invariant level of $X_{ic}$ and the firm's baseline count for
clause type $c$), and $\alpha_{ft}$ is a flexible time effect. We report
three choices for $\alpha_{ft}$: a year fixed effect, a clause type
$\times$ year fixed effect (which absorbs clause-specific secular trends),
and a mode-union $\times$ year fixed effect (which absorbs union-specific
time trends). Standard errors are clustered at the firm level.

The treatment indicator is $T_t = \mathbf{1}\{cba\_period_{t} \geq 3\}$,
i.e.\ post-Súmula renegotiation. Identification of $\beta$ comes from
within firm $\times$ clause-type variation: for a given $(i,c)$ cell, how
much the post-period count moves relative to that same firm's other
clause-type cells where the firm did or did not have a pre-treatment
mismatch.

\paragraph{Predictions.} Under the content-diffusion hypothesis,
%
\begin{itemize}
  \item $\beta_{\text{gap}} > 0$: treated firms add clauses where their
    connected partners had more.
  \item $\beta_{\text{surplus}} \approx 0$: partners cannot teach a firm
    about clauses they themselves are short of, so a positive level of
    \emph{own} pre-treatment surplus should not predict post growth. A
    significantly negative $\beta_{\text{surplus}}$ instead points to
    \emph{convergence} (firms shed clauses they had above the local mean)
    or to mechanical reallocation under a quasi-fixed renegotiation budget.
  \item $\beta_{\text{raw\_gap}}$ is the OLS slope under the linearity
    restriction $\beta_{\text{gap}} = -\beta_{\text{surplus}}$. We
    additionally estimate equation~\eqref{eq:mechanism_spec} with both
    $\text{gap}_{ic} \times T_t$ and $\text{surplus}_{ic} \times T_t$
    entered jointly and report a Wald test of this linearity restriction.
\end{itemize}

\paragraph{Placebo.} As a pre-trend check, we re-estimate
equation~\eqref{eq:mechanism_spec} on the pre-period subset
$cba\_period \in \{1, 2\}$, replacing $T_t$ with
$T^{\text{plac}}_t = \mathbf{1}\{cba\_period_{t} = 2\}$. The same firm
$\times$ clause-type cells are compared between their first and second
pre-Súmula periods. A coefficient of zero is consistent with no
differential pre-trend in clause counts on either side of the mismatch.
A non-zero placebo coefficient would indicate that the mismatch measures
already correlated with clause-count growth before Súmula 277, and would
require attribution of any post-period coefficient to a continuing trend
rather than the reform itself.

\paragraph{Mirror to the firm-level similarity exercise.} The mechanism
test is the treated-side mirror of the untreated-firm CBA-similarity
exercise (Section~\ref{sec:cba_similarity}, table~\ref{tab:cba_similarity}).
There, more-connected untreated firms have their full CBA vectors compared
to a connected-treated reference. Here we ask the same question
clause-by-clause from the treated side: do treated firms re-template
specifically toward what their connected partners carry? Both reduce to
the same content-diffusion null, evaluated with different aggregations
and from opposite sides of the worker flow.
